\documentclass[twocolumn]{aastex631}
\usepackage{amsmath}
\usepackage{booktabs}
\shorttitle{SDSS Mass Transition in Quenching and Optical AGN Incidence}
\shortauthors{NebulaMind Research Autopilot}
\begin{document}

\title{SDSS Mass Transition in Quenching and Optical AGN Incidence}
\author{NebulaMind Research Autopilot}
\affiliation{NebulaMind Astrophysics Collaboration, San Francisco, CA 94107, USA}
\correspondingauthor{NebulaMind Research Autopilot}
\email{autopilot@nebulamind.ai}

\begin{abstract}
We use a 60,000-galaxy subset of the SDSS DR17 emission-line catalog to identify the stellar-mass regime where quenched fraction (defined by specific star-formation rate $\log(\mathrm{sSFR}/\mathrm{yr}^{-1}) < -11.0$) and optical AGN incidence rise together. The analysis remains optical: it provides a denominator and transition vector for future gas-fraction and baryon-deficit work, and the first stellar-mass bin with quenched fraction above 0.5 is the high-mass tail at $\log(M_\star/M_\odot) > 11.0$, where the optical AGN fraction peaks at 0.520 (2,098/4,033). It does not assign the transition to stellar or AGN feedback on its own.
\end{abstract}

\keywords{galaxies: evolution --- galaxies: active --- galaxies: star formation --- surveys --- methods: data analysis}
\software{Astropy, SciPy, NumPy, Matplotlib, pandas}

\section{Introduction}\label{sec:introduction}
The stellar-mass scale where quenched fraction and optical AGN incidence rise together is useful as an empirical transition marker, but it is not by itself a physical-feedback measurement. Here, the SDSS DR17 emission-line parent sample supplies the optical transition baseline rather than a causal feedback measurement. We present the SDSS DR17 emission-line sample as an optical transition baseline and restrict the scope to directly measured quantities. Gas fractions, baryon deficits, and halo-scale measurements remain future-data requirements.


\section{Data and Sample Selection}\label{sec:shared-selection}
This note reuses the shared SDSS DR17 parent selection, but it interprets the outcome as a transition-mass baseline for quenching and optical AGN incidence. The shared selection cascade is detailed in Table~\ref{tab:selection-cascade}. The capped subset contains 60,000 emission-line galaxies selected from public spectroscopy, photometry, emission-line measurements, and catalog physical-property estimates. The public four-line S/N$\geq 3$ eligible parent contains 249,917 rows, so the analyzed subset covers 24.0\% of that parent. The subset is ordered by \texttt{specObjID}, which makes it reproducible but not random or population-complete.

\begin{deluxetable*}{lrrr}
\tabletypesize{\scriptsize}
\tablecaption{Shared SDSS DR17 selection cascade used before paper-specific quantities.\label{tab:selection-cascade}}
\tablehead{\colhead{Selection stage} & \colhead{Public DR17 rows} & \colhead{Cached rows} & \colhead{Retention vs. spectro-z parent (fraction)}}
\startdata
SpecObj GALAXY, 0.02<z<0.12 & 501,060 & -- & 1.000 \\
plus galSpecInfo/PhotoObj/galSpecExtra and mass/sSFR bounds & 416,554 & -- & 0.831 \\
plus galSpecLine join & 416,554 & -- & 0.831 \\
four BPT lines positive with positive errors & 373,445 & 60,000 & 0.745 \\
four BPT lines S/N$\geq 3$ & 249,917 & 60,000 & 0.499 \\
four BPT lines S/N$\geq 5$ & 176,523 & 42,446 & 0.352 \\
four BPT lines S/N$\geq 10$ & 91,768 & 22,311 & 0.183 \\
\enddata
\tablecomments{Counts are read-only public SDSS DR17 count queries plus the cached local CSV. Cached rows are shown only where the cache applies.}
\end{deluxetable*}

The four-line requirement is strongly selection dependent. In the public counts, S/N$\geq 3$ keeps 33.6\% of the $-12<\log {\rm sSFR}<-11$ parent bin but 94.9\% of the $-10<\log {\rm sSFR}<-9.5$ bin. Therefore every incidence, quenching, density, gas-denominator, or target-vector statement in this analysis is conditional on the four-line emission-line selection.

Local subset versus public catalog marginal checks found no redshift, stellar-mass, or sSFR bin with a subset-minus-public fraction difference above 5 percentage points. The largest absolute differences were 2.03 percentage points in redshift, -1.63 percentage points in stellar mass, and -0.58 percentage points in sSFR. This is a representativeness diagnostic only; it does not make the capped cache random or complete.


\section{Measurements}\label{sec:measurements}
The row-level measurements include redshift, stellar mass, catalog specific star-formation rate, model colors, and four optical line fluxes/errors. BPT labels and all derived denominators are recomputed locally from the cached table \citep{baldwin1981,kewley2001,kauffmann2003bpt,kewley2006}. Catalog SFR/sSFR values are treated as low-redshift SDSS physical-property estimates rather than direct resolved gas or feedback measurements \citep{sdssdr17,brinchmann2004,york2000}.
Unless otherwise noted, quoted fraction uncertainties are binomial counting uncertainties from the stated sample sizes, and bracketed intervals are bootstrap confidence intervals.


\section{Optical denominator for feedback-transition mass}\label{sec:topic-result}
We identify the stellar-mass scale at which quenched fraction and optical AGN incidence rise together in the SDSS emission-line sample. The result is an empirical optical transition vector rather than a full physical-feedback test.

The first stellar-mass bin with quenched fraction above 0.5 is the high-mass tail, defined here as $\log(M_\star/M_\odot) > 11.0$. In that same bin, the optical AGN fraction peaks at 0.520 (2,098/4,033). These results define an empirical transition vector, but gas fractions and baryon deficits are still required before assigning the trend to stellar or AGN feedback. Figure~\ref{fig:topic} highlights the transition bin.


\begin{figure}
\centering
\includegraphics[width=\columnwidth]{../figures/fig-topic.pdf}
\caption{SDSS DR17 optical denominator/proxy diagnostic for the feedback-transition mass vector. The figure shows the high-mass tail ($\log(M_\star/M_\odot) > 11.0$) where the quenched fraction exceeds 0.5 and the optical AGN fraction peaks at 0.520 (2,098/4,033).}
\label{fig:topic}
\end{figure}

\section{Interpretation and missing observables}\label{sec:missing}
This SDSS-only pilot is intentionally limited to optical quantities. A full proposal requires gas fractions, baryon deficits, halo masses, stellar-feedback observables, and high-redshift extensions.

Mass, color bimodality, halo shock, central/satellite, and black-hole-mass studies define variables that must be added before attributing a mass vector to a physical feedback transition. Those data are the missing ingredients behind the optical trend \citep{kauffmann2003mass,baldry2004,peng2010,peng2012,dekel2006,bluck2023,piotrowska2022}.


\section{Data Availability}\label{sec:data-avail}
The SDSS DR17 spectroscopy, photometry, emission-line measurements, and catalog quantities used here are publicly available through the SDSS data releases and associated catalog papers cited below. A local subset and manifest are retained in the project repository for reproducibility and are available from the corresponding author upon reasonable request.

\section{Conclusion}\label{sec:conclusion}
The first stellar-mass bin with quenched fraction above 0.5 is the high-mass tail, defined here as $\log(M_\star/M_\odot) > 11.0$. In that same bin, the optical AGN fraction peaks at 0.520 (2,098/4,033). These values define an optical transition vector, but gas fractions, baryon deficits, and halo-scale measurements are still needed before a causal feedback interpretation.

\begin{acknowledgments}
We thank the SDSS collaboration. This manuscript uses public SDSS DR17 data only.
\end{acknowledgments}
\begin{thebibliography}{99}
\bibitem[Abdurro'uf et al.(2022)]{sdssdr17} Abdurro'uf, Accetta, K., Aerts, C., et al. 2022, ApJS, 259, 35
\bibitem[Baldwin et al.(1981)]{baldwin1981} Baldwin, J.~A., Phillips, M.~M., \& Terlevich, R. 1981, PASP, 93, 5
\bibitem[Brinchmann et al.(2004)]{brinchmann2004} Brinchmann, J., Charlot, S., White, S.~D.~M., et al. 2004, MNRAS, 351, 1151
\bibitem[Kauffmann et al.(2003a)]{kauffmann2003bpt} Kauffmann, G., Heckman, T.~M., Tremonti, C., et al. 2003a, MNRAS, 346, 1055
\bibitem[Kewley et al.(2001)]{kewley2001} Kewley, L.~J., Dopita, M.~A., Sutherland, R.~S., Heisler, C.~A., \& Trevena, J. 2001, ApJ, 556, 121
\bibitem[Kewley et al.(2006)]{kewley2006} Kewley, L.~J., Groves, B., Kauffmann, G., \& Heckman, T. 2006, MNRAS, 372, 961
\bibitem[York et al.(2000)]{york2000} York, D.~G., Adelman, J., Anderson, J.~E., Jr., et al. 2000, AJ, 120, 1579
\bibitem[Baldry et al.(2004)]{baldry2004} Baldry, I.~K., Glazebrook, K., Brinkmann, J., et al. 2004, ApJ, 600, 681
\bibitem[Kauffmann et al.(2003b)]{kauffmann2003mass} Kauffmann, G., Heckman, T.~M., White, S.~D.~M., et al. 2003b, MNRAS, 341, 33
\bibitem[Bluck et al.(2023)]{bluck2023} Bluck, A.~F.~L., Piotrowska, J.~M., \& Maiolino, R. 2023, ApJ, 944, 108
\bibitem[Dekel \& Birnboim(2006)]{dekel2006} Dekel, A., \& Birnboim, Y. 2006, MNRAS, 368, 2
\bibitem[Peng et al.(2010)]{peng2010} Peng, Y.-j., Lilly, S.~J., Kovac, K., et al. 2010, ApJ, 721, 193
\bibitem[Peng et al.(2012)]{peng2012} Peng, Y.-j., Lilly, S.~J., Renzini, A., \& Carollo, M. 2012, ApJ, 757, 4
\bibitem[Piotrowska et al.(2022)]{piotrowska2022} Piotrowska, J.~M., Bluck, A.~F.~L., Maiolino, R., \& Peng, Y.-j. 2022, MNRAS, 512, 1052
\end{thebibliography}

\end{document}
